% =====================================================================
%  Section I. Introduction
% =====================================================================
\section{INTRODUCTION}\label{sec:intro}

The success of the quark-constituent picture both for resonances and for
deep-inelastic electron and neutrino processes makes it difficult to
believe quarks do not exist.  The problem is that quarks have not been
seen.  This suggests that quarks, for some reason, cannot appear as
separate particles in a final state.  A number of speculations have been
offered as to how this might happen.${}^{1}$

Independently of the quark problem, Schwinger observed many years ago${}^{2}$
that the vector mesons of a gauge theory can have a nonzero mass if
vacuum polarization totally screens the charges in a gauge theory.
Schwinger illustrated this result with the exact solution of quantum
electrodynamics in one space and one time dimension, where the photon
acquires a mass $\sim e^{2}$ for any nonzero charge $e$ [$e$ has
dimensions of $(\mathrm{mass})^{1/2}$ in this theory].  Schwinger
suggested that the same effect could occur in four dimensions for
sufficiently large couplings.

Further study of the Schwinger model by Lowenstein and Swieca${}^{3}$
and Casher, Kogut, and Susskind${}^{4}$ has shown that the asymptotic
states of the model contain only massive photons, not electrons.
Nevertheless, as Casher \emph{et al.} have shown in detail, the
electrons are present in deep-inelastic processes and behave like free
pointlike particles over short times and short distances.  The
polarization effects which prevent the appearance of electrons in the
final state take place on a longer time scale (longer than $1/m_{\gamma}$,
where $m_{\gamma}$ is the photon mass).

A new mechanism which keeps quarks bound will be proposed in this paper.
The mechanism applies to gauge theories only.  The mechanism will be
illustrated using the strong-coupling limit of a gauge theory in
four-dimensional space-time.  However, the model discussed here has a
built-in ultraviolet cutoff, and in the strong-coupling limit all
particle masses (including the gauge-field masses) are much larger than
the cutoff; in consequence the theory is far from covariant.

The confinement mechanism proposed here is soft (long-time scale).
However, in the model discussed here the cutoff spoils the possibility
of free pointlike behavior for the quarks.

The model discussed in this paper is a gauge theory set up on a
four-dimensional Euclidean lattice.  The inverse of the lattice spacing
$a$ serves as an ultraviolet cutoff.  The use of a Euclidean space
(i.e., imaginary instead of real times) instead of a Lorentz space is
not a serious restriction; the energy eigenstates (including scattering
states) of the lattice theory can be determined from the
``transfer-matrix'' formalism as has been discussed by Suri${}^{5}$ and
reviewed by Wilson and Kogut.${}^{6}$  A brief discussion of the
transfer-matrix method is given in Sec.~III.

In Schwinger's speculations about four dimensions, the photon mass
would be zero for any charge $e$ less than a critical coupling $e_{c}$;
for $e > e_{c}$, the photon mass would be nonzero and vary with $e$.
Figure~\ref{fig:1} shows how a plot of $m_{\gamma}$ vs $e$ might look.
The point $e_{c}$ is a point of nonanalyticity.  Similar nonanalytic
points, called critical points, occur in solid-state physics at certain
types of phase transitions.  Consider, for instance, a ferromagnet in
the absence of an external field.  For any temperature above the Curie
temperature $T_{c}$, the spontaneous magnetization $M$ is 0.  Below
$T_{c}$, $M$ is nonzero and a function of temperature.  At $T_{c}$ there
may be either a first-order phase transition (in which case $M$ is
discontinuous at $T_{c}$) or a second-order phase transition (critical
point) for which $M \to 0$ as $T \to T_{c}$ from either side.

\begin{figure}[t]
  \centering
  \includegraphics[width=0.42\textwidth]{fig1}
  \refstepcounter{figure}%
  \caption*{Fig. 1. Speculative plot of photon mass vs renormalized
    charge $e$, in unknown units. The transition at $e_{c}$ is
    second-order (see text).}
  \label{fig:1}
\end{figure}

By analogy with the solid-state situation one can think of the
transition from zero to nonzero photon mass as a change of phase: this
analogy is best understood by imagining the particles of quantum
electrodynamics to be the excitations of a medium (the ether).  In this
case it is the ether which undergoes a change of phase at $e_{c}$.  There
is again a question whether this change of phase is first-order (cf.\
Fig.~\ref{fig:2}) or second-order (Fig.~\ref{fig:1}).  (Coleman and
Weinberg${}^{7}$ have found a nontrivial example of a first-order
transition in another context.)

The model discussed in this paper is a single Abelian gauge field
coupled (with strength $g$) to massive quarks.  In weak coupling the
gauge field behaves like a normal free zero-mass field (despite
modifications introduced in the lattice quantization) and the quarks are
unbound.  In strong coupling the gauge field is massive and the quarks
are bound, showing the existence of the second phase.  Thus there should
be a phase transition at some intermediate value of $g$.  Nothing is
known about this transition at the present time.

\begin{figure}[t]
  \centering
  \includegraphics[width=0.42\textwidth]{fig2}
  \refstepcounter{figure}%
  \caption*{Fig. 2. Speculative plot of photon mass vs renormalized
    charge $e$ if there is a first-order transition at $e_{c}$.}
  \label{fig:2}
\end{figure}

The quantization procedure and strong-coupling approximation described
in this paper can be applied to non-Abelian gauge theories also.  This
will be explained briefly in Sec.~III.

An extraordinary feature of the strong-coupling expansion of the lattice
theory (see Sec.~IV) is that it has the same general structure as the
relativistic string models of hadrons.${}^{8}$  The vacuum expectation
values of the gauge theory involve (in the strong-coupling expansion)
sums over all quark paths and sums over all surfaces connecting these
paths; the surfaces are generated by the gauge field treated in strong
coupling.  The paths and surfaces are defined on a discrete lattice.
There are geometrical difficulties in relating the surfaces on the
lattice to the continuum surfaces of the string models; it is not known
at present whether these difficulties can be overcome.

In Sec.~II the nature of the quark binding mechanism will be discussed,
qualitatively.  In Sec.~III the gauge theory will be formulated on a
discrete lattice, both classically and quantum mechanically.  In
Sec.~IV the strong-coupling expansion for the lattice gauge theory is
explained.  In Sec.~V a cursory discussion of weak coupling is given.
In Sec.~VI there is a brief discussion of the problem of Lorentz
invariance and the relation to string models.
